% GENERATED FILE. Edit canonical lessons, then run npm run generate:books.
% canonical-source-hash: fnv1a64:67dbedc129e36066
% canonical-lessons: HI-C462-abhiman, HI-C462-itrana, HI-C462-sekhi, HI-C462-akarna, HI-C462-nakhra

\chapter{Showing Off}
\label{ch:hi-showing-off}
\hlchaptermodality{462}

\begin{chapteropening}
\textbf{By the end of this chapter:} \emph{I can say that someone is showing off.}

Complete the last lesson of chapter 462: I can say that someone is showing off.
\end{chapteropening}

\section[\texorpdfstring{\hi{अभिमान} (abhimān)}{अभिमान (abhimān)}]{\hi{अभिमान} (abhimān) — self-esteem; pride}

\label{lesson:HI-C462-abhiman}



\begin{quote}
Before the new one: say the Hindi for to blush, then the Hindi for brazen.
\end{quote}

\subsection*{You'll want to know: \hi{अभिमान}}
\textbf{\hi{अभिमान}} — \emph{abhimān} — \textquotedblleft{}self-esteem; pride\textquotedblright{}.

\begin{quote}
\hi{अरुण} \hi{को} \hi{अपने} \hi{काम} \hi{पर} \hi{अभिमान} \hi{है।}
\end{quote}

\begin{quote}
\emph{Aruṇ ko apne kām par abhimān hai.}
\end{quote}

\begin{quote}
Arun takes pride in his work.
\end{quote}

Masculine.

\begin{cousinweb}
\textbf{\hi{अभिमान}} (\emph{abhimān}) is a tatsama, borrowed from Sanskrit \emph{abhimāna}, \textquotedblleft{}a high opinion of oneself, pride\textquotedblright{} (Monier-Williams). It can be honest pride or vanity: the sentence decides.
\end{cousinweb}

\begin{grammarlens}[title={What you've built}]
One more word for describing what you lived and felt.
\end{grammarlens}

\subsection*{Your turn}
\begin{itemize}
\raggedright
  \item \emph{Say it:} \emph{abhimān}
  \item \emph{Say it:} \emph{abhimān}, once more
  \item \emph{Say it:} \emph{Aruṇ ko apne kām par abhimān hai.} — Arun takes pride in his work.
  \item \emph{Recall:} say the Hindi for to blush, then the Hindi for brazen, then say \emph{abhimān} again
\end{itemize}

\begin{tcolorbox}[breakable,colback=teal!4,colframe=teal!35!black,title={Before you move on}]
What does \hi{अभिमान} mean? (\textquotedblleft{}Self-esteem; pride\textquotedblright{}.) Say it once more.
\end{tcolorbox}
\section[\texorpdfstring{\hi{इतराना} (itrānā)}{इतराना (itrānā)}]{\hi{इतराना} (itrānā) — to put on airs}

\label{lesson:HI-C462-itrana}



\begin{quote}
Before the new one: say the Hindi for brazen, then the Hindi for self-esteem.
\end{quote}

\subsection*{You'll want to know: \hi{इतराना}}
\textbf{\hi{इतराना}} — \emph{itrānā} — \textquotedblleft{}to put on airs\textquotedblright{}.

\begin{quote}
\hi{नए} \hi{कपड़े} \hi{पहनकर} \hi{मीरा} \hi{की} \hi{बहन} \hi{इतरा} \hi{रही} \hi{थी।}
\end{quote}

\begin{quote}
\emph{nae kapṛe pahankar Mīrā kī bahan itrā rahī thī.}
\end{quote}

\begin{quote}
Wearing her new clothes, Mira's sister was putting on airs.
\end{quote}

Past (she …): \hi{वह} \hi{इतराई}, \emph{vah itrāī}.

\begin{cousinweb}
\textbf{\hi{इतराना}} (\emph{itrānā}) is inherited: CDIAL lists Hindi \emph{itrānā}, \textquotedblleft{}to strut\textquotedblright{}, under Sanskrit \emph{itvará}, \textquotedblleft{}going, walking\textquotedblright{} (CDIAL 1566).
\end{cousinweb}

\begin{grammarlens}[title={What you've built}]
One more word for describing what you lived and felt.
\end{grammarlens}

\subsection*{Your turn}
\begin{itemize}
\raggedright
  \item \emph{Say it:} \emph{itrānā}
  \item \emph{Say it:} \emph{itrānā}, once more
  \item \emph{Say it:} \emph{nae kapṛe pahankar Mīrā kī bahan itrā rahī thī.} — Wearing her new clothes, Mira's sister was putting on airs.
  \item \emph{Recall:} say the Hindi for brazen, then the Hindi for self-esteem, then say \emph{itrānā} again
\end{itemize}

\begin{tcolorbox}[breakable,colback=teal!4,colframe=teal!35!black,title={Before you move on}]
What does \hi{इतराना} mean? (\textquotedblleft{}To put on airs\textquotedblright{}.) Say it once more.
\end{tcolorbox}
\section[\texorpdfstring{\hi{शेख़ी} (śekhī)}{शेख़ी (śekhī)}]{\hi{शेख़ी} (śekhī) — a boast; bragging}

\label{lesson:HI-C462-sekhi}



\begin{quote}
Before the new one: say the Hindi for self-esteem, then the Hindi for to put on airs.
\end{quote}

\subsection*{You'll want to know: \hi{शेख़ी}}
\textbf{\hi{शेख़ी}} — \emph{śekhī} — \textquotedblleft{}a boast; bragging\textquotedblright{}.

\begin{quote}
\hi{अरुण} \hi{का} \hi{दोस्त} \hi{हमेशा} \hi{अपनी} \hi{गाड़ी} \hi{की} \hi{शेख़ी} \hi{मारता} \hi{है।}
\end{quote}

\begin{quote}
\emph{Aruṇ kā dost hameśā apnī gāṛī kī śekhī mārtā hai.}
\end{quote}

\begin{quote}
Arun's friend is always boasting about his car.
\end{quote}

Feminine.

\begin{cousinweb}
\textbf{\hi{शेख़ी}} (\emph{śekhī}) is usually traced to Persian \emph{shēkhī}, \textquotedblleft{}the airs of a sheikh\textquotedblright{}, hence boasting. \emph{śekhī mārnā} is \textquotedblleft{}to brag\textquotedblright{}.
\end{cousinweb}

\begin{grammarlens}[title={What you've built}]
One more word for describing what you lived and felt.
\end{grammarlens}

\subsection*{Your turn}
\begin{itemize}
\raggedright
  \item \emph{Say it:} \emph{śekhī}
  \item \emph{Say it:} \emph{śekhī}, once more
  \item \emph{Say it:} \emph{Aruṇ kā dost hameśā apnī gāṛī kī śekhī mārtā hai.} — Arun's friend is always boasting about his car.
  \item \emph{Recall:} say the Hindi for self-esteem, then the Hindi for to put on airs, then say \emph{śekhī} again
\end{itemize}

\begin{tcolorbox}[breakable,colback=teal!4,colframe=teal!35!black,title={Before you move on}]
What does \hi{शेख़ी} mean? (\textquotedblleft{}A boast; bragging\textquotedblright{}.) Say it once more.
\end{tcolorbox}
\section[\texorpdfstring{\hi{अकड़ना} (akaṛnā)}{अकड़ना (akaṛnā)}]{\hi{अकड़ना} (akaṛnā) — to strut; to stiffen}

\label{lesson:HI-C462-akarna}



\begin{quote}
Before the new one: say the Hindi for to put on airs, then the Hindi for a boast.
\end{quote}

\subsection*{You'll want to know: \hi{अकड़ना}}
\textbf{\hi{अकड़ना}} — \emph{akaṛnā} — \textquotedblleft{}to strut; to stiffen\textquotedblright{}.

\begin{quote}
\hi{जीत} \hi{के} \hi{बाद} \hi{खिलाड़ी} \hi{बहुत} \hi{अकड़} \hi{रहा} \hi{था।}
\end{quote}

\begin{quote}
\emph{jīt ke bād khilāṛī bahut akaṛ rahā thā.}
\end{quote}

\begin{quote}
After the win, the player was strutting about.
\end{quote}

Past (he …): \hi{वह} \hi{अकड़} \hi{गया}, \emph{vah akaṛ gayā}.

\begin{cousinweb}
\textbf{\hi{अकड़ना}} (\emph{akaṛnā}) is a tadbhava: CDIAL lists Hindi \emph{akaṛnā}, \textquotedblleft{}to be numbed\textquotedblright{}, under \emph{*ākkaḍa}, a form it reconstructs (CDIAL 1013). Stiff limbs and a stiff, haughty back are both \emph{akaṛnā}.
\end{cousinweb}

\begin{grammarlens}[title={What you've built}]
One more word for describing what you lived and felt.
\end{grammarlens}

\subsection*{Your turn}
\begin{itemize}
\raggedright
  \item \emph{Say it:} \emph{akaṛnā}
  \item \emph{Say it:} \emph{akaṛnā}, once more
  \item \emph{Say it:} \emph{jīt ke bād khilāṛī bahut akaṛ rahā thā.} — After the win, the player was strutting about.
  \item \emph{Recall:} say the Hindi for to put on airs, then the Hindi for a boast, then say \emph{akaṛnā} again
\end{itemize}

\begin{tcolorbox}[breakable,colback=teal!4,colframe=teal!35!black,title={Before you move on}]
What does \hi{अकड़ना} mean? (\textquotedblleft{}To strut; to stiffen\textquotedblright{}.) Say it once more.
\end{tcolorbox}
\section[\texorpdfstring{\hi{नख़रा} (nakhrā)}{नख़रा (nakhrā)}]{\hi{नख़रा} (nakhrā) — fussiness; airs and graces}

\label{lesson:HI-C462-nakhra}



\begin{quote}
Before the new one: say the Hindi for a boast, then the Hindi for to strut.
\end{quote}

\subsection*{You'll want to know: \hi{नख़रा}}
\textbf{\hi{नख़रा}} — \emph{nakhrā} — \textquotedblleft{}fussiness; airs and graces\textquotedblright{}.

\begin{quote}
\hi{मेरी} \hi{बहन} \hi{खाने} \hi{में} \hi{बहुत} \hi{नख़रे} \hi{करती} \hi{है।}
\end{quote}

\begin{quote}
\emph{merī bahan khāne meṁ bahut nakhre kartī hai.}
\end{quote}

\begin{quote}
My sister is very fussy about food.
\end{quote}

Masculine.

\begin{cousinweb}
\textbf{\hi{नख़रा}} (\emph{nakhrā}) is usually traced to Persian \emph{nakhra}, \textquotedblleft{}coquetry, airs\textquotedblright{}. \emph{nakhre karnā} is to make a fuss.
\end{cousinweb}

\begin{grammarlens}[title={What you've built}]
One more word for describing what you lived and felt.
\end{grammarlens}

\subsection*{Your turn}
\begin{itemize}
\raggedright
  \item \emph{Say it:} \emph{nakhrā}
  \item \emph{Say it:} \emph{nakhrā}, once more
  \item \emph{Say it:} \emph{merī bahan khāne meṁ bahut nakhre kartī hai.} — My sister is very fussy about food.
  \item \emph{Recall:} say the Hindi for a boast, then the Hindi for to strut, then say \emph{nakhrā} again
\end{itemize}

\begin{tcolorbox}[breakable,colback=teal!4,colframe=teal!35!black,title={Before you move on}]
What does \hi{नख़रा} mean? (\textquotedblleft{}Fussiness; airs and graces\textquotedblright{}.) Say it once more.
\end{tcolorbox}
